% =====================================================================
% PARTE III — Lavorare nel sistema
% =====================================================================
\part{Lavorare nel sistema}

\noindent\itshape Quattro capitoli che chiudono la traversata. Prima un
preambolo (\S\ref{cap:prima}) per dichiarare cosa cambia entrando in
Parte~III: i sorgenti Agda useranno pattern matching, e questo ha un
costo che va esplicitato (quinto cero). Poi la logica non importata ma
riconosciuta dentro F-I-E-C, fino a osservare Curry-Howard come gesto
finale (\S\ref{cap:logica}). Un kit di costruzioni canoniche --- i
\emph{tipi notevoli} (\S\ref{cap:notevoli}) --- che fa da palestra
prima delle prove vere. E infine l'esercizio di Stump sulla
commutativit\`a di \texttt{\&\&} letto come tre F-I-E-C intrecciati
(\S\ref{cap:and-comm}), apertura del compendio al capitolo~1 del libro
in oggetto di studio.\upshape

\bigskip

% ---------------------------------------------------------------------
\chapter{Prima del compendio}\label{cap:prima}

\noindent {\spacedlowsmallcaps{Cosa cambia, entrando in Parte~III}}. In Parte~II si scrivevano regole. Da qui in poi si \emph{leggono} sorgenti, ed \`e onesto dichiarare in apertura che cosa si sta facendo quando una clausola Agda viene tradotta in F-I-E-C.

\bigskip

\noindent\textbf{La situazione.} \emph{Verified Functional Programming in Agda} (e qualsiasi sorgente Agda reale) usa il \emph{pattern matching} come modo standard di definire funzioni e dimostrazioni. Una clausola del tipo

\begin{agda}
f : (b : 𝔹) → C b
f tt = ...
f ff = ...
\end{agda}

non scrive $\indB$ da nessuna parte. Il nostro compito di lettura sar\`a riconoscere, sotto ogni clausola di questo genere, il F-I-E-C in atto: che cosa \`e l'eliminazione, che cosa la computazione, che cosa l'introduzione del risultato.

Nel caso semplice di $\mathbb{B}$ la traduzione \`e meccanica. Su $\mathbb{N}$ e $\mathrm{Vec}$ comincer\`a a non esserlo: il pattern matching dipendente \emph{eccede} ci\`o che $\indN$, $\indV$ ecc.\ riescono a fare da soli. Va detto chiaramente, ed \`e il contenuto del quinto cero.

\section*{Il quinto cero}

\begin{cero}[Il pattern matching dipendente eccede MLTT pura]\label{cero:pm}
Il pattern matching che vediamo in Agda non \`e zucchero per gli eliminatori di Parte~II. Nella sua forma generale o \textbf{assume l'assioma K} (\emph{uniqueness of identity proofs}, ogni prova di $x \identita x$ \`e definitoriamente $\refl$), oppure --- in modalit\`a \texttt{--without-K} --- si appoggia a un'\textbf{elaborazione non banale} (Goguen-McBride-McKinna, Cockx-Abel) basata su unificazione di indici, \emph{dot patterns}, \emph{absurd patterns}, \emph{with-abstraction}.

In entrambi i casi: ci\`o che scriviamo come clausole con costruttori a sinistra \emph{non \`e} immediatamente un'applicazione di $\indB$, $\indN$, $\Jelim$ ecc. C'\`e un compilatore in mezzo. Quando in questo libro leggiamo una clausola come ``tre F-I-E-C intrecciati'' (\S\ref{cap:and-comm}) stiamo facendo un'operazione di traduzione che, fuori dai casi pi\`u semplici, non \`e meccanica.
\end{cero}

\section*{Dove morde la profondit\`a: un esempio breve}

L'esempio canonico \`e il matching di $\refl$ contro una prova generica $p : x \identita y$. Vedremo la situazione completa in \S\ref{cap:notevoli} con $\mathsf{sym}$ e $\mathsf{trans}$, ma il punto \`e fissabile in poche righe.

Scrivere

\begin{agda}
sym : {A : Set} {x y : A} → x ≡ y → y ≡ x
sym refl = refl
\end{agda}

\noindent \emph{sembra} dire: ``c'\`e un solo caso, quello in cui la prova \`e $\refl$, e nel ramo $\refl$ il sistema sa che $y = x$''. Ma da dove viene quel ``il sistema sa che $y = x$''? In MLTT pura, l'unico modo di estrarre informazione da $p : x \identita y$ \`e $\Jelim$ (\S\ref{cap:identita}): si specifica il motivo $C : (y : A) \to (x \identita y) \to \Type$ e si abita il caso $C\,x\,\refl$.

\medskip

\textbf{Il matching di Agda fa di pi\`u}: nel ramo $\refl$ \emph{istanzia} $y := x$ nel resto del contesto. Quell'istanziazione, in piena generalit\`a, \`e l'assioma K. In \texttt{--without-K} si controlla che l'istanziazione sia legittima sul piano dell'unificazione --- e questa verifica \`e un algoritmo, non una regola di MLTT.

\begin{oss}[Conseguenza pratica]
Su tipi senza uguaglianze proposizionali non banali fra costruttori --- $\mathbb{B}$, $\mathbb{N}$ per gli esempi facili --- la traduzione clausola$\to$eliminatore \`e meccanica e non chiama in causa K. \texttt{\&\&-comm} (\S\ref{cap:and-comm}) \`e proprio uno di questi casi. Quando lavoreremo su prove che \emph{matchano} su $\refl$ in modo essenziale (\S\ref{cap:notevoli} per $\mathsf{sym}$, $\mathsf{trans}$, $\mathsf{cong}$, $\mathsf{subst}$) il cero~\ref{cero:pm} sar\`a in scena, e va saputo.
\end{oss}

\section*{Patto di lettura per Parte~III}

\noindent\textbf{(i)} Ogni clausola Agda di Parte~III verr\`a accompagnata, almeno la prima volta che il pattern compare, dalla sua traduzione esplicita in F-I-E-C (eliminatore + motivo + casi base/passo).

\noindent\textbf{(ii)} Quando la traduzione \`e \emph{meccanica} (nessun matching essenziale su $\refl$, nessun uso di dot pattern o absurd pattern non banale) ce lo segneremo come ``elaborazione triviale''.

\noindent\textbf{(iii)} Quando \`e \emph{non meccanica} il cero~\ref{cero:pm} sar\`a in scena, lo dichiareremo, e indicheremo a quale lato (K o \texttt{--without-K}) ci stiamo appoggiando.

\bigskip

\noindent Con questo patto in tasca, possiamo entrare in Parte~III sapendo cosa stiamo importando e cosa no.

% ---------------------------------------------------------------------
\chapter{La logica che affiora}\label{cap:logica}

\noindent {\spacedlowsmallcaps{Nessun cero nuovo qui}}. Sono \textbf{osservazioni}, non assunzioni. La logica non \`e importata: \`e gi\`a l\`i nella struttura F-I-E-C.

\bigskip

\noindent\textbf{Gesto del capitolo.} Non diremo ``ora decidiamo che i tipi sono proposizioni e i termini sono prove''. Quella sarebbe un'assunzione --- un cero. Diremo invece: \emph{guardiamo le regole F-I-E-C dei tipi che abbiamo costruito in Parte~II, e notiamo che hanno la forma delle regole della logica intuizionista del primo ordine}. La differenza fra le due mosse \`e tutto il capitolo. La prima impone una lettura; la seconda riconosce una struttura. Il nome di questa coincidenza --- \emph{corrispondenza di Curry-Howard} --- arriva alla fine come osservazione su quanto avr\`a gi\`a parlato da s\'e.

\section*{Implicazione: $\Pi$ non dipendente}

Quando $B$ non dipende da $x$, la regola di formazione di $\Pi$ (\S\ref{cap:pi}) si specializza:
\[
\inferrule{A : \Type \\ B : \Type}{A \to B : \Type}
\]

\noindent\textbf{Lettura logica.} Premesse: $A$ \`e una proposizione, $B$ \`e una proposizione. Conclusione: $A \to B$ \`e una proposizione --- l'implicazione $A \Rightarrow B$.

Le altre due regole di $\Pi$ si leggono altrettanto direttamente:
\begin{itemize}
\item Introduzione $\lambda x.\,b : A \to B$ sotto l'ipotesi $x : A$: \textbf{dimostrazione per ipotesi}. Si assume $A$, si deriva $B$, si chiude scaricando l'ipotesi. \`E la regola $\Rightarrow$-introduzione della deduzione naturale.
\item Eliminazione $f\,a : B$ dati $f : A \to B$ e $a : A$: \textbf{modus ponens}. La regola $\Rightarrow$-eliminazione.
\end{itemize}

\noindent Niente \`e stato deciso. Lo scaricamento delle ipotesi --- gesto centrale della deduzione naturale --- \emph{coincide con} l'introduzione di $\Pi$, che era stata scritta per ragioni di calcolo (descrivere come si costruisce una funzione).

\section*{Congiunzione: $\Sigma$ non dipendente}

Quando $B$ non dipende da $x$, la formazione di $\Sigma$ (\S\ref{cap:sigma}) d\`a:
\[
\inferrule{A : \Type \\ B : \Type}{A \times B : \Type}
\]

\noindent\textbf{Lettura logica.} $A \times B$ \`e la proposizione $A \wedge B$.

\begin{itemize}
\item Introduzione $(a, b) : A \times B$ da $a : A$ e $b : B$: \textbf{regola $\wedge$-introduzione}.
\item Eliminazioni $\mathsf{fst}$, $\mathsf{snd}$ (caso particolare di $\indS$ con motivo costante): \textbf{regole $\wedge$-eliminazione 1 e 2}.
\end{itemize}

\section*{Falsit\`a: $\bot$}

Il tipo vuoto in Agda:

\begin{agda}
data ⊥ : Set where
  -- nessun costruttore
\end{agda}

Le sue quattro regole F-I-E-C esplicitate:

\[
\inferrule{ }{\bot : \Type}
\qquad
\textit{(nessuna I)}
\qquad
\inferrule{e : \bot \\ C : \bot \to \Type}{\mathsf{ind}\text{-}\bot(e, C) : C\,e}
\qquad
\textit{(nessuna C)}
\]

Una Formazione, \emph{zero} Introduzioni, un'Eliminazione, \emph{zero} Computazioni. Eliminatore di $\bot$ in Agda:

\begin{agda}
⊥-elim : {A : Set} → ⊥ → A
⊥-elim ()
\end{agda}

L'\emph{absurd pattern} \texttt{()} \`e la realizzazione sintattica di ``zero costruttori, niente da matchare, l'esaustivit\`a si chiude da sola''.

\noindent\textbf{Lettura logica.}
\begin{itemize}
\item $\bot$ non \`e abitato (nessuna introduzione): \textbf{$\bot$ \`e la falsit\`a}.
\item L'eliminatore d\`a $\bot \to A$ per ogni $A$: \textbf{ex falso quodlibet}.
\end{itemize}

\section*{Negazione: una definizione, non una primitiva}

Non c'\`e una regola F-I-E-C per la negazione. \`E una \emph{definizione}:
\[
\neg A \;\;\coloneqq\;\; A \to \bot
\]

\noindent ``Negare $A$'' significa ``mostrare che da $A$ si deriverebbe l'assurdo''. Lettura intuizionista canonica della negazione, ottenuta gratis dall'avere $\to$ e $\bot$. Una prova di $\neg A$ \`e una funzione che riceve un ipotetico abitante di $A$ e produce un abitante di $\bot$ (cosa impossibile in concreto: \`e la dimostrazione di impossibilit\`a).

\section*{Quantificatore universale: $\Pi$ dipendente}

Senza specializzare, la regola di formazione di $\Pi$ (\S\ref{cap:pi}) \`e gi\`a:
\[
\inferrule{A : \Type \\ x : A \vdash B\,x : \Type}{(x : A) \to B\,x : \Type}
\]

\noindent\textbf{Lettura logica.} ``Per ogni $x$ di tipo $A$, vale $B\,x$''. \`E $\forall x : A.\,B\,x$.

L'introduzione $\lambda x.\,b\,x$ \`e: \emph{costruisci la prova di $B\,x$ uniformemente in $x$, e ottieni la prova del $\forall$}. L'eliminazione $f\,a : B\,a$ \`e: \emph{specializza il $\forall$ all'elemento $a$}.

\noindent Stessa regola di $\Pi$, due letture. Non ne abbiamo aggiunta nessuna nuova.

\section*{Quantificatore esistenziale: $\Sigma$ dipendente}

Stesso gesto sull'altra regola (\S\ref{cap:sigma}):
\[
\inferrule{A : \Type \\ x : A \vdash B\,x : \Type}{\Sigma\,x \in A,\,B\,x : \Type}
\]

\noindent\textbf{Lettura logica.} ``Esiste $x$ di tipo $A$ tale che $B\,x$''. \`E $\exists x : A.\,B\,x$.

L'introduzione $(a, b) : \Sigma\,x \in A,\,B\,x$ richiede un \emph{testimone} $a$ e una \emph{prova} $b : B\,a$. Questa \`e una propriet\`a forte della logica intuizionista: l'esistenza \`e \textbf{costruttiva}, una prova di $\exists x.\,P\,x$ \emph{esibisce} l'$x$ in questione. Non si pu\`o dimostrare ``esiste qualcosa con propriet\`a $P$'' senza dire \emph{cosa} ce l'ha.

\begin{oss}[Confronto con la matematica classica]
In matematica classica si dimostra spesso ``esiste $x$ tale che $\ldots$'' assumendo per assurdo la negazione e derivando una contraddizione. La prova esiste, il testimone no. In MLTT questa strategia \emph{non funziona}: per abitare $\Sigma\,x \in A,\,B\,x$ devi produrre $(a, b)$. \`E lo stesso fenomeno che vedremo a fine capitolo per la doppia negazione e LEM.
\end{oss}

\section*{Disgiunzione: il coproduct $A \uplus B$}\label{sec:coproduct}

Qui c'\`e una connettiva che ci manca. Non l'abbiamo introdotta in Parte~II perch\'e non serviva per la traversata dei sei tipi fondativi. La introduciamo adesso, perch\'e la lettura logica la chiama: \emph{disgiunzione costruttiva}.

\medskip

\noindent F-I-E-C completa del coproduct (o \emph{somma disgiunta}):

\paragraph{F --- Formazione.}
\[
\inferrule{A : \Type \\ B : \Type}{A \uplus B : \Type}
\]

\paragraph{I --- Introduzione.} (Due costruttori.)
\[
\inferrule{a : A}{\mathsf{inj}_1\,a : A \uplus B}
\qquad
\inferrule{b : B}{\mathsf{inj}_2\,b : A \uplus B}
\]

\paragraph{E --- Eliminazione.}
\[
\inferrule{
e : A \uplus B \\
C : (A \uplus B) \to \Type \\
f : (a : A) \to C\,(\mathsf{inj}_1\,a) \\
g : (b : B) \to C\,(\mathsf{inj}_2\,b)
}{
\mathsf{ind}\text{-}\uplus(e, C, f, g) : C\,e
}
\]

\paragraph{C --- Computazione.}
\begin{align*}
\mathsf{ind}\text{-}\uplus(\mathsf{inj}_1\,a, C, f, g) &\definitorio f\,a \\
\mathsf{ind}\text{-}\uplus(\mathsf{inj}_2\,b, C, f, g) &\definitorio g\,b
\end{align*}

Stesso schema visto in Parte~II: due costruttori, nessuna ricorsione (perch\'e n\'e $\mathsf{inj}_1$ n\'e $\mathsf{inj}_2$ contengono istanze di $A \uplus B$), eliminazione per casi.

In Agda:

\begin{agda}
data _⊎_ (A B : Set) : Set where
  inj₁ : A → A ⊎ B
  inj₂ : B → A ⊎ B
\end{agda}

\noindent\textbf{Lettura logica.} $A \uplus B$ \`e la proposizione $A \vee B$.

\begin{itemize}
\item Introduzioni: regole $\vee$-introduzione 1 e 2.
\item Eliminazione: regola di \emph{ragionamento per casi} ($\vee$-eliminazione). Per concludere $C$ da $A \vee B$, mostra come concludere $C$ assumendo $A$, e come concludere $C$ assumendo $B$.
\end{itemize}

\textbf{Costruttivit\`a della disgiunzione.} Una prova di $A \uplus B$ ti dice \emph{quale dei due} vale. Non c'\`e modo di abitarlo senza decidere fra $\mathsf{inj}_1$ e $\mathsf{inj}_2$. \`E lo stesso fenomeno dell'esistenziale: la disgiunzione \`e \emph{informativa}.

\section*{Uguaglianza}

Il tipo identit\`a $\_\identita\_$ (\S\ref{cap:identita}) \`e gi\`a la lettura logica dell'uguaglianza fra termini di uno stesso tipo. $\refl$ \`e l'unica introduzione, e via $\Jelim$ si dimostrano le propriet\`a usuali (riflessivit\`a per costruzione, simmetria, transitivit\`a, sostituzione --- le faremo in \S\ref{cap:notevoli}). Niente da aggiungere qui, salvo notare che \textbf{non} \`e una connettiva: \`e una famiglia di proposizioni indicizzata da una coppia di termini.

\section*{Cosa \emph{non} c'\`e: la dimensione costruttiva}

Tutto quello che abbiamo letto come logica \`e \textbf{intuizionista}. Significa, in concreto:

\begin{itemize}
\item \textbf{Terzo escluso non derivabile.} Il tipo $A \uplus \neg A$, in generale, \emph{non} \`e abitato da niente che si possa costruire uniformemente in $A$. Sarebbe una funzione $(A : \Type) \to A \uplus \neg A$ --- non c'\`e modo di scriverla in MLTT senza un postulato esterno. Logicamente: $A \vee \neg A$ non \`e un teorema dello schema F-I-E-C.
\item \textbf{Doppia negazione non eliminabile.} $\neg\neg A \to A$, in generale, non \`e abitato. Una prova di $\neg\neg A$ \`e una funzione $(A \to \bot) \to \bot$: ti dice che ``non si pu\`o produrre una funzione $A \to \bot$''. Da questa informazione \emph{non} si estrae un abitante di $A$. La differenza ``non c'\`e contro-esempio'' / ``c'\`e un esempio'' \`e costitutiva.
\end{itemize}

\noindent\textbf{Osservazione, non difetto.} La logica che affiora qui non \`e una logica \emph{mancante di qualcosa}: \`e una logica \emph{diversa}, dove la nozione di prova \`e pi\`u stringente. Una prova di $A$ \emph{esibisce} un abitante; una prova di $\exists x.\,P\,x$ \emph{esibisce} il testimone; una prova di $A \vee B$ \emph{decide} quale dei due vale. Per recuperare la matematica classica si pu\`o postulare LEM o DNE come ceri aggiuntivi (Stump lo fa esplicitamente in alcuni contesti); ma quello \`e un gesto separato, e va segnato come tale.

\section*{Tabella riassuntiva}

\begin{center}
\renewcommand{\arraystretch}{1.2}
\begin{tabular}{lll}
\toprule
\textit{Logica intuizionista} & \textit{Tipo} & \textit{Da quale F-I-E-C} \\
\midrule
$A \Rightarrow B$       & $A \to B$              & $\Pi$ non dip.\ (\S\ref{cap:pi}) \\
$A \wedge B$            & $A \times B$           & $\Sigma$ non dip.\ (\S\ref{cap:sigma}) \\
$\bot$ (falsit\`a)      & $\bot$                 & questa sezione, $\bot$ \\
$\neg A$                & $A \to \bot$           & definizione, $\to$ + $\bot$ \\
$\forall x : A.\,P\,x$  & $(x : A) \to P\,x$     & $\Pi$ dip.\ (\S\ref{cap:pi}) \\
$\exists x : A.\,P\,x$  & $\Sigma\,x \in A,\,P\,x$ & $\Sigma$ dip.\ (\S\ref{cap:sigma}) \\
$A \vee B$              & $A \uplus B$           & questa sezione, \S\ref{sec:coproduct} \\
$a = b$ in $A$          & $a \identita b$        & tipo identit\`a (\S\ref{cap:identita}) \\
\midrule
\multicolumn{3}{c}{\textit{non derivabili in generale}} \\
\midrule
LEM: $A \vee \neg A$    & $A \uplus (A \to \bot)$ & --- \\
DNE: $\neg\neg A \Rightarrow A$ & $((A \to \bot) \to \bot) \to A$ & --- \\
\bottomrule
\end{tabular}
\end{center}

\section*{Curry-Howard, come osservazione}

\noindent {\spacedlowsmallcaps{Quello che abbiamo appena descritto ha un nome}}. La coincidenza fra struttura dei tipi e struttura della logica intuizionista del primo ordine si chiama \textbf{corrispondenza di Curry-Howard} (talvolta \emph{Curry-Howard-de Bruijn}, o letta come ``isomorfismo proposizioni-come-tipi''). La presentano cos\`i:

\medskip
\begin{center}
proposizioni $\;\longleftrightarrow\;$ tipi \qquad prove $\;\longleftrightarrow\;$ termini
\end{center}
\medskip

\noindent Nei testi di logica dei tipi questa \`e quasi sempre l'apertura: ``decidiamo che tipi sono proposizioni, e dimostrare significa abitare''. Qui no. \`E arrivata alla fine, come constatazione di una corrispondenza gi\`a \emph{verificata caso per caso} nelle sezioni precedenti, e non come postulato di partenza.

\medskip

\noindent\textbf{Perch\'e la differenza importa.} Dichiarare la corrispondenza in apertura sarebbe un cero: ``assumiamo che tipi e proposizioni siano la stessa cosa''. Riconoscerla in chiusura \`e un'osservazione: ``avendo costruito i tipi per ragioni di calcolo (descrivere funzioni dipendenti, induzione, coppie), ci troviamo davanti che \emph{erano gi\`a} le regole della logica intuizionista''. Il nome resta utile come riferimento bibliografico; il \emph{contenuto} \`e tutto ci\`o che abbiamo gi\`a fatto.

\begin{oss}[Cero che non c'\`e]
La corrispondenza di Curry-Howard non figura nell'elenco dei ceri (\S\ref{cap:ceri}). Non perch\'e ce ne siamo dimenticati: perch\'e non \`e un'assunzione di MLTT. \`E una propriet\`a osservabile delle regole che MLTT ha gi\`a per altri motivi.
\end{oss}

% ---------------------------------------------------------------------
\chapter{Tipi notevoli}\label{cap:notevoli}

\epigraph{Come i limiti notevoli in analisi: poche costruzioni canoniche, da imparare una volta, da riconoscere ovunque.}{}

\noindent {\spacedlowsmallcaps{Un kit di attrezzi}}. Prima di entrare nella prova reale di Stump (\S\ref{cap:and-comm}), raccogliamo le costruzioni che si incontrano dappertutto nelle prove di Verified Functional Programming in Agda, scritte due volte: una alla Little Typer (con eliminatore esplicito), una in Agda (con pattern matching). \`E la palestra dove fissare l'\emph{esercizio centrale} di \S\ref{sec:stile}: leggere F-I-E-C sotto ogni clausola.

\bigskip

\noindent Cinque sezioni, in ordine di profondit\`a crescente: funzioni booleane (\S\ref{sec:not-funz-bool}), funzioni numeriche (\S\ref{sec:funz-nat}), propriet\`a del tipo identit\`a (\S\ref{sec:notevoli-id}), tipi piccoli di servizio (\S\ref{sec:tipi-piccoli}), operazioni su $\mathrm{Vec}$ (\S\ref{sec:notevoli-vec}). L'ultima sezione \`e il primo posto dove il cero~\ref{cero:pm} morde davvero.

\section{Funzioni booleane}\label{sec:not-funz-bool}

\subsection*{$\mathsf{not}$}

\paragraph{Definizione (Little Typer).} Eliminatore $\indB$ con motivo costante:
\[
\mathsf{not}\,b \;\definitorio\; \indB(b,\;\lambda\_.\,\mathbb{B},\;\ff,\;\tto)
\]

Letto: ``per costruire un $\mathbb{B}$ a partire da $b$, da' $\ff$ se $b = \tto$, $\tto$ se $b = \ff$''.

\paragraph{Definizione (Agda).}
\begin{agda}
not : 𝔹 → 𝔹
not tt = ff
not ff = tt
\end{agda}

Le due clausole \emph{sono} i due argomenti $c_t, c_f$ di $\indB$. Elaborazione triviale: niente matching su $\refl$, niente dot pattern.

\subsection*{$\_\&\&\_$ e $\_||\_$}

\paragraph{Definizione (Little Typer).}
\begin{align*}
b_1 \andd b_2 &\;\definitorio\; \indB(b_1,\;\lambda\_.\,\mathbb{B},\;b_2,\;\ff) \\
b_1 \,\|\, b_2 &\;\definitorio\; \indB(b_1,\;\lambda\_.\,\mathbb{B},\;\tto,\;b_2)
\end{align*}

Curry sul primo argomento: $b_2$ \`e nel contesto quando $\indB$ entra in scena, quindi sopravvive ai due rami.

\paragraph{Definizione (Agda).}
\begin{agda}
_&&_ : 𝔹 → 𝔹 → 𝔹
tt && b = b
ff && b = ff

_||_ : 𝔹 → 𝔹 → 𝔹
tt || b = tt
ff || b = b
\end{agda}

\begin{oss}[Cucitura con \S\ref{cap:and-comm}]
La definizione di \texttt{\_\&\&\_} appena vista \`e \textbf{esattamente} la stessa che useremo nel capitolo successivo per dimostrare \texttt{\&\&-comm}. Le ``regole di computazione di \texttt{\&\&}'' di cui parler\`a \S\ref{cap:and-comm} sono i due passi $\beta$ di $\indB$ applicato a $\tto$ e $\ff$.
\end{oss}

\subsection*{$\mathsf{if\_then\_else\_}$}

\paragraph{Definizione (Little Typer).} Per qualsiasi tipo $A$:
\[
\mathsf{if}\,b\,\mathsf{then}\,x\,\mathsf{else}\,y \;\definitorio\; \indB(b,\;\lambda\_.\,A,\;x,\;y)
\]

\paragraph{Definizione (Agda).}
\begin{agda}
if_then_else_ : {A : Set} → 𝔹 → A → A → A
if tt then x else y = x
if ff then x else y = y
\end{agda}

\noindent\textbf{Da fissare.} Le quattro funzioni booleane di base sono \emph{quattro istanze} dello stesso eliminatore $\indB$, distinte solo dal motivo e dai due rami. Nessuna nuova regola.

\section{Funzioni numeriche}\label{sec:funz-nat}

Stesso gesto su $\mathbb{N}$, ma ora i passi induttivi ricevono l'ipotesi induttiva.

\subsection*{$\_+\_$}

\paragraph{Definizione (Little Typer).} Induzione sul primo argomento, motivo $\mathbb{N} \to \mathbb{N}$ (cio\`e: ``$\mathsf{add}\,n$ \`e una funzione $\mathbb{N} \to \mathbb{N}$''):
\[
\mathsf{add}\,n \;\definitorio\; \indN(n,\;\lambda\_.\,\mathbb{N}\to\mathbb{N},\;\lambda m.\,m,\;\lambda k.\,\lambda \mathit{rec}.\,\lambda m.\,\suc\,(\mathit{rec}\,m))
\]

\noindent\textbf{Lettura dei pezzi:}
\begin{itemize}
\item motivo $\lambda\_.\mathbb{N}\to\mathbb{N}$: per ogni $n$, $\mathsf{add}\,n$ \`e una funzione $\mathbb{N}\to\mathbb{N}$;
\item caso $\zero$: $\lambda m.\,m$ --- aggiungere zero \`e l'identit\`a;
\item passo $\suc$: ricevo $k$, l'ipotesi induttiva $\mathit{rec} : \mathbb{N}\to\mathbb{N}$ (cio\`e $\mathsf{add}\,k$), e produco $\lambda m.\,\suc\,(\mathit{rec}\,m)$.
\end{itemize}

\paragraph{Definizione (Agda).}
\begin{agda}
_+_ : ℕ → ℕ → ℕ
zero  + m = m
suc k + m = suc (k + m)
\end{agda}

La ``chiamata ricorsiva'' \texttt{k + m} nel secondo ramo \`e \textbf{esattamente} il $\mathit{rec}\,m$ del passo induttivo di sopra. Il typechecker accetta perch\'e la ricorsione \`e su \texttt{k}, strutturalmente pi\`u piccolo di \texttt{suc k} (\S\ref{cap:nat}).

\subsection*{$\_*\_$}

\paragraph{Definizione (Little Typer).}
\[
\mathsf{mul}\,n \;\definitorio\; \indN(n,\;\lambda\_.\,\mathbb{N}\to\mathbb{N},\;\lambda m.\,\zero,\;\lambda k.\,\lambda \mathit{rec}.\,\lambda m.\,m + \mathit{rec}\,m)
\]

\paragraph{Definizione (Agda).}
\begin{agda}
_*_ : ℕ → ℕ → ℕ
zero  * m = zero
suc k * m = m + (k * m)
\end{agda}

\begin{oss}[Pattern: ``funzione binaria via induzione sul primo argomento'']
Sia \texttt{\_+\_} sia \texttt{\_*\_} usano lo stesso schema: $\indN$ sul primo argomento, motivo $\mathbb{N}\to\mathbb{N}$, caso base = funzione su $m$, passo = funzione su $m$ che usa la ricorsione. La maggior parte delle funzioni $\mathbb{N}\to\mathbb{N}\to X$ in Stump cap.~2 segue questo schema. Riconoscerlo basta a leggere mezzo capitolo del libro.
\end{oss}

\section{Propriet\`a del tipo identit\`a}\label{sec:notevoli-id}

\noindent {\spacedlowsmallcaps{Qui entra in scena il cero~\ref{cero:pm}}}. Tutte e quattro le costruzioni che seguono sono casi particolari di $\Jelim$ (\S\ref{cap:identita}), e tutte e quattro le loro versioni Agda fanno matching su $\refl$.

\subsection*{$\mathsf{sym}$ --- simmetria}

\paragraph{Definizione (Little Typer).} Dato $p : x \identita y$, vogliamo $y \identita x$. Si applica $\Jelim$ con motivo $C\,y\,p \definitorio y \identita x$. Caso base ($p = \refl$, $y = x$): $C\,x\,\refl = x \identita x$, abitato da $\refl$.
\[
\mathsf{sym}\,p \;\definitorio\; \Jelim(p,\;\lambda y\,p.\,y \identita x,\;\refl)
\]

\paragraph{Definizione (Agda).}
\begin{agda}
sym : {A : Set} {x y : A} → x ≡ y → y ≡ x
sym refl = refl
\end{agda}

\noindent\textbf{Cosa fa il matching di $\refl$.} Quando Agda vede la clausola \texttt{sym refl = ...}, istanzia $y := x$ nel resto del contesto. Quell'istanziazione \`e ci\`o che il cero~\ref{cero:pm} chiama in causa: in MLTT pura l'unico modo di ``sapere che $y = x$'' \`e specificare il motivo di $\Jelim$ e abitarlo nel caso $\refl$. Qui Agda lo fa da s\'e.

\subsection*{$\mathsf{trans}$ --- transitivit\`a}

\paragraph{Definizione (Little Typer).} Dati $p : x \identita y$, $q : y \identita z$, vogliamo $x \identita z$. Si applica $\Jelim$ su $p$ con motivo $C\,y\,p \definitorio y \identita z \to x \identita z$. Caso base: $C\,x\,\refl = x \identita z \to x \identita z$, abitato dall'identit\`a $\lambda q.\,q$.
\[
\mathsf{trans}\,p\,q \;\definitorio\; \Jelim(p,\;\lambda y\,p.\,y \identita z \to x \identita z,\;\lambda q.\,q)\,q
\]

\paragraph{Definizione (Agda).}
\begin{agda}
trans : {A : Set} {x y z : A} → x ≡ y → y ≡ z → x ≡ z
trans refl q = q
\end{agda}

Una sola clausola: il matching su $p = \refl$ istanzia $y := x$, e il tipo da abitare diventa $x \identita z$, che \`e gi\`a $q$.

\subsection*{$\mathsf{cong}$ --- congruenza}

\paragraph{Definizione (Little Typer).} Data $f : A \to B$ e $p : x \identita y$, vogliamo $f\,x \identita f\,y$. Motivo: $C\,y\,p \definitorio f\,x \identita f\,y$. Caso base: $C\,x\,\refl = f\,x \identita f\,x$, abitato da $\refl$.
\[
\mathsf{cong}\,f\,p \;\definitorio\; \Jelim(p,\;\lambda y\,p.\,f\,x \identita f\,y,\;\refl)
\]

\paragraph{Definizione (Agda).}
\begin{agda}
cong : {A B : Set} {x y : A} (f : A → B) → x ≡ y → f x ≡ f y
cong f refl = refl
\end{agda}

\subsection*{$\mathsf{subst}$ --- trasporto}

\paragraph{Definizione (Little Typer).} Data una propriet\`a $P : A \to \Type$, una prova $p : x \identita y$, e $\mathit{px} : P\,x$, vogliamo $P\,y$. Motivo: $C\,y\,p \definitorio P\,x \to P\,y$. Caso base: $C\,x\,\refl = P\,x \to P\,x$, abitato da $\lambda\mathit{px}.\,\mathit{px}$.
\[
\mathsf{subst}\,P\,p\,\mathit{px} \;\definitorio\; \Jelim(p,\;\lambda y\,p.\,P\,x \to P\,y,\;\lambda \mathit{px}.\,\mathit{px})\,\mathit{px}
\]

\paragraph{Definizione (Agda).}
\begin{agda}
subst : {A : Set} (P : A → Set) {x y : A} → x ≡ y → P x → P y
subst P refl px = px
\end{agda}

\begin{oss}[Quattro variazioni sullo stesso tema]
$\mathsf{sym}$, $\mathsf{trans}$, $\mathsf{cong}$, $\mathsf{subst}$ sono \textbf{quattro istanze di $\Jelim$}, distinte solo dal motivo. Il fatto che in Agda si scrivano tutte come una o due clausole con $\refl$ a sinistra \`e l'effetto del cero~\ref{cero:pm}: l'elaborazione del pattern matching trova il motivo da sola. Non magia, lavoro di unificazione.
\end{oss}

\section{Tipi piccoli di servizio}\label{sec:tipi-piccoli}

\subsection*{$\top$ --- il tipo unitario}

F-I-E-C completa (un costruttore, nessuna ricorsione):

\paragraph{F.} $\top : \Type$.

\paragraph{I.} $\star : \top$ (un solo costruttore senza premesse).

\paragraph{E.}
\[
\inferrule{u : \top \\ C : \top \to \Type \\ c : C\,\star}{\mathsf{ind}\text{-}\top(u, C, c) : C\,u}
\]

\paragraph{C.} $\mathsf{ind}\text{-}\top(\star, C, c) \definitorio c$.

\medskip

In Agda:

\begin{agda}
data ⊤ : Set where
  ★ : ⊤
\end{agda}

\noindent\textbf{Usi.} $\top$ \`e il tipo che ``non porta informazione'': abita la posizione di \emph{vero} nella logica, e serve come ``riempitivo'' quando un motivo deve restituire qualcosa di banale (lo vedremo immediato in \S\ref{sec:notevoli-vec} per $\mathsf{head}$).

\subsection*{$\neg$ --- negazione (richiamo)}

Gi\`a vista in \S\ref{cap:logica}: $\neg A \coloneqq A \to \bot$. Niente F-I-E-C nuova, \`e una definizione. Vale la pena ricordarla qui perch\'e compare in pratica sempre quando si lavora con $\mathsf{Dec}$, che chiameremo se serve:
\[
\mathsf{Dec}\,A \;\coloneqq\; A \uplus \neg A
\]

\noindent ``$A$ \`e decidibile'' = c'\`e o una prova di $A$ o una refutazione di $A$, e si sa quale.

\section{Operazioni su $\mathrm{Vec}\,A\,n$}\label{sec:notevoli-vec}

\noindent {\spacedlowsmallcaps{Dove il pattern matching comincia a non essere zucchero}}. Le funzioni che seguono sono ancora gestibili a mano con $\indV$ (\S\ref{cap:vec}), ma una di esse --- $\mathsf{head}$ --- richiede un trucco sul motivo che il pattern matching nasconde con grazia. \`E un buon posto per vedere il cero~\ref{cero:pm} al lavoro su tipi induttivi non banali.

\subsection*{$\mathsf{head}$ e $\mathsf{tail}$ --- e perch\'e il caso $\nil$ non c'\`e}

\paragraph{Definizione (Agda).}
\begin{agda}
head : {A : Set} {n : ℕ} → Vec A (suc n) → A
head (cons a as) = a

tail : {A : Set} {n : ℕ} → Vec A (suc n) → Vec A n
tail (cons a as) = as
\end{agda}

\noindent\textbf{La cosa che salta all'occhio.} Manca il caso $\nil$. Eppure Agda accetta come esaustivo. Perch\'e?

Perch\'e l'indice $\suc\,n$ nel tipo dell'argomento \emph{esclude} $\nil$: $\nil$ ha tipo $\mathrm{Vec}\,A\,\zero$, non $\mathrm{Vec}\,A\,(\suc\,n)$ per alcun $n$. L'unificatore di Agda riconosce questo e impone esaustivit\`a sui soli costruttori compatibili con l'indice. Questo \`e \textbf{esattamente} il tipo di lavoro che il cero~\ref{cero:pm} chiama in causa.

\paragraph{Versione Little Typer.} Volendo scrivere $\mathsf{head}$ con $\indV$ puro, l'unico modo \`e definire una funzione \emph{totale} su $\mathrm{Vec}\,A\,n$ qualunque $n$ sia, scegliendo un motivo che restituisca $\top$ nel caso $\zero$ e $A$ nel caso $\suc$:
\[
C : (n : \mathbb{N}) \to \mathrm{Vec}\,A\,n \to \Type
\qquad
C\,\zero\,\_ \definitorio \top
\qquad
C\,(\suc\,k)\,\_ \definitorio A
\]

Allora:
\[
\mathsf{head}'\,vs \;\definitorio\; \indV(vs,\;C,\;\star,\;\lambda k.\,\lambda a.\,\lambda \mathit{as}.\,\lambda\mathit{rec}.\,a)
\]

Per ottenere $\mathsf{head}$ ``vera'' --- quella che pretende un vettore di lunghezza $\suc\,n$ --- si specializza $\mathsf{head}'$ al caso in cui l'indice \`e $\suc\,n$, dove il motivo riduce ad $A$:
\[
\mathsf{head}\,(vs : \mathrm{Vec}\,A\,(\suc\,n)) \;\definitorio\; \mathsf{head}'\,vs : A
\]

\begin{oss}[Cosa abbiamo guadagnato col pattern matching]
La versione Agda \texttt{head (cons a as) = a} \`e \emph{una riga}. La versione Little Typer richiede di inventare un motivo che usa $\top$ come ``segnaposto'' per il caso che non interessa, definire una funzione totale, e poi specializzarla. Il contenuto matematico \`e lo stesso, ma la versione Agda \`e \emph{infinitamente} pi\`u leggibile. \`E il guadagno reale del cero~\ref{cero:pm}: il prezzo (assumere K, o appoggiarsi all'unificazione di \texttt{--without-K}) compra leggibilit\`a sostanziale.
\end{oss}

\subsection*{$\_\mathbin{+\!\!+}\_$ --- concatenazione}

\paragraph{Definizione (Agda).}
\begin{agda}
_++_ : {A : Set} {m n : ℕ} → Vec A m → Vec A n → Vec A (m + n)
nil         ++ ys = ys
(cons a as) ++ ys = cons a (as ++ ys)
\end{agda}

\paragraph{Definizione (Little Typer).} $\indV$ sul primo vettore, motivo che parametrizza sul vettore di destinazione:
\[
C : (m : \mathbb{N}) \to \mathrm{Vec}\,A\,m \to \Type
\qquad
C\,m\,\_ \definitorio \mathrm{Vec}\,A\,n \to \mathrm{Vec}\,A\,(m + n)
\]

\noindent\textbf{Cosa sta lavorando.} Nel caso $\nil$ ($m = \zero$): $C\,\zero\,\nil = \mathrm{Vec}\,A\,n \to \mathrm{Vec}\,A\,(\zero + n)$. Per la \emph{regola di computazione di $\_+\_$} (\S\ref{sec:funz-nat}), $\zero + n \definitorio n$, quindi il tipo \`e $\mathrm{Vec}\,A\,n \to \mathrm{Vec}\,A\,n$ --- abitato da $\lambda ys.\,ys$. Nel caso $\cons\,a\,\mathit{as}$ ($m = \suc\,k$): tipo $\mathrm{Vec}\,A\,n \to \mathrm{Vec}\,A\,(\suc\,k + n) = \mathrm{Vec}\,A\,(\suc\,(k+n))$, abitato da $\lambda ys.\,\cons\,a\,(\mathit{rec}\,ys)$.

\medskip

\noindent\textbf{Da fissare.} La regola di computazione di $\_+\_$ (definitoria, vista in \S\ref{sec:funz-nat}) lavora \emph{dentro il tipo}. Il typechecker la usa silenziosamente per verificare che $\mathrm{Vec}\,A\,(\zero + n)$ e $\mathrm{Vec}\,A\,n$ siano lo stesso tipo. \`E il Cero~\ref{cero:dipendenza} (un termine appare dentro un tipo) in azione operativa.

\subsection*{$\mathsf{map}$}

\paragraph{Definizione (Agda).}
\begin{agda}
map : {A B : Set} {n : ℕ} → (A → B) → Vec A n → Vec B n
map f nil         = nil
map f (cons a as) = cons (f a) (map f as)
\end{agda}

\paragraph{Definizione (Little Typer).} $\indV$ con motivo $C\,n\,\_ \definitorio \mathrm{Vec}\,B\,n$ (non dipende dal vettore, solo dalla lunghezza, perch\'e $\mathsf{map}$ preserva la lunghezza per costruzione):
\[
\mathsf{map}\,f\,vs \;\definitorio\; \indV(vs,\;\lambda n\,\_.\,\mathrm{Vec}\,B\,n,\;\nil,\;\lambda k.\,\lambda a.\,\lambda \mathit{as}.\,\lambda\mathit{rec}.\,\cons\,(f\,a)\,\mathit{rec})
\]

\begin{oss}[Il motivo costante sul vettore]
Nei tre notevoli di Vec di sopra, il motivo dipende dalla \emph{lunghezza} ma quasi mai dal \emph{vettore} stesso. Questo \`e tipico delle funzioni che ``lavorano in parallelo'' (mappa, zip, ecc.) rispetto a quelle che ``leggono'' il vettore (head, tail, lookup). Vederlo nei motivi rende esplicito cosa una funzione fa.
\end{oss}

\section{Cosa portare nel capitolo successivo}

In ordine di necessit\`a per \S\ref{cap:and-comm}:

\begin{itemize}
\item La definizione di \texttt{\_\&\&\_} come $\indB$ sul primo argomento con motivo costante (\S\ref{sec:not-funz-bool}). \`E ci\`o che render\`a meccaniche le ``riduzioni di \texttt{\&\&}'' nel quadro F-I-E-C di \texttt{\&\&-comm}.
\item L'introduzione di $\refl$ del tipo identit\`a, $\Jelim$ in standby, e il fatto che nel capitolo successivo \emph{non} useremo $\mathsf{sym}/\mathsf{trans}/\mathsf{cong}/\mathsf{subst}$ --- la prova si chiude solo con $\refl$, senza matching essenziale su prove di uguaglianza. Significa che il cero~\ref{cero:pm} \emph{non} \`e in scena nel caso \texttt{\&\&-comm}; arriver\`a nei capitoli successivi quando entrer\`a $\mathbb{N}$.
\end{itemize}

% ---------------------------------------------------------------------
\chapter{Anatomia di una prova: \texttt{\&\&-comm} in Iowa Agda Library}\label{cap:and-comm}

\noindent {\spacedlowsmallcaps{Mettiamo a frutto tutto il libro per leggere una prova reale}}. Ogni cosa che succede qui poggia su F-I-E-C gi\`a visti. \texttt{\_\&\&\_} la conosciamo dal kit dei notevoli (\S\ref{sec:not-funz-bool}); $\refl$ e il tipo identit\`a da \S\ref{cap:identita}; l'eliminatore $\indB$ da \S\ref{cap:bool}. E il \emph{patto di lettura} stabilito in \S\ref{cap:prima}: ogni clausola Agda verr\`a tradotta in F-I-E-C, segnalando se la traduzione \`e meccanica o se chiama in causa il cero~\ref{cero:pm}. \textbf{Anticipazione, per questo capitolo}: \`e meccanica. Il cero~\ref{cero:pm} non entra in scena --- entrer\`a quando faremo prove su $\mathbb{N}$ (cap.~2 di Stump).

\bigskip

L'esercizio del capitolo~1 di \emph{Verified Functional Programming in Agda} (Stump, libreria IAL):

\begin{agda}
module ch01.Booleans where
open import bool         -- definisce 𝔹, _&&_
open import bool-thms    -- teoremi precostruiti
open import eq           -- definisce _≡_ e refl (è il tipo identità di sezione 11)

&&-comm : ∀ (b1 b2 : 𝔹) → b1 && b2 ≡ b2 && b1
&&-comm tt tt = refl
&&-comm tt ff = refl
&&-comm ff tt = refl
&&-comm ff ff = refl
\end{agda}

Lo srotoliamo in tre tempi: che cos'\`e il \textbf{tipo}, che cos'\`e la \textbf{definizione di \texttt{\&\&}}, che cosa fa il \textbf{pattern matching}.

\section{Il tipo, strato per strato}

Il tipo \`e puro zucchero per due $\Pi$ annidati. Tolto lo zucchero:

\begin{agda}
-- strato 3 (file originale): forall + parentesi raggruppate
&&-comm : ∀ (b1 b2 : 𝔹) → b1 && b2 ≡ b2 && b1

-- strato 2: forall e' zucchero per Pi con tipo annotato (sez. forall)
&&-comm :   (b1 b2 : 𝔹) → b1 && b2 ≡ b2 && b1

-- strato 1: parentesi raggruppate e' zucchero per Pi annidati (sez. pi-azione)
&&-comm :   (b1 : 𝔹) → ((b2 : 𝔹) → (b1 && b2) ≡ (b2 && b1))

-- strato 0: MLTT puro, due applicazioni della regola di formazione di Pi (sez. pi)
--   con il tipo del risultato finale che e' un tipo identita' (sez. identita)
\end{agda}

Quindi \texttt{\&\&-comm} deve abitare \textbf{due $\Pi$ annidati}. Curry dipendente (\S\ref{cap:pi-azione}): nel secondo $\Pi$, \texttt{b1} \`e gi\`a nel contesto. Il tipo del risultato finale, \texttt{b1 \&\& b2 $\equiv$ b2 \&\& b1}, \`e il \textbf{tipo identit\`a} del linguaggio (\S\ref{cap:identita}) --- non l'uguaglianza definitoria.

\section{La definizione di \texttt{\&\&} come regole di computazione}

In IAL \texttt{\_\&\&\_} \`e definita:

\begin{agda}
_&&_ : 𝔹 → 𝔹 → 𝔹
tt && b = b
ff && b = ff
\end{agda}

\noindent\`E lo stesso \texttt{\_\&\&\_} di \S\ref{sec:not-funz-bool}: $\indB$ sul primo argomento, motivo costante $\lambda\_.\mathbb{B}$, rami $b_2$ e $\ff$. Le due clausole \emph{sono regole di computazione} sotto mentite spoglie. Cio\`e, il typechecker riduce automaticamente:
\begin{align*}
\tto \andd b &\definitorio b \\
\ff \andd b &\definitorio \ff
\end{align*}

\begin{oss}[Punto cruciale]
\texttt{\&\&} matcha \emph{solo sul primo argomento}. Se il primo \`e una variabile non istanziata, non riduce. Quindi:

\smallskip
\begin{itemize}
\item $\tto \andd b_2$ riduce a $b_2$. \checkmark
\item $\ff \andd b_2$ riduce a $\ff$. \checkmark
\item $b_1 \andd b_2$ con $b_1$ variabile aperta: \textbf{non riduce}. $\times$
\item $b_1 \andd \tto$ con $b_1$ variabile aperta: \textbf{non riduce}. $\times$
\end{itemize}
\end{oss}

\section{Il pattern matching come tre cose insieme}\label{sec:tre-fiec}

Le quattro clausole \texttt{\&\&-comm tt tt = refl} ecc.\ nascondono \textbf{tre F-I-E-C distinti} che lavorano in sequenza.

\paragraph{(a) Eliminazione di $\mathbb{B}$ (\S\ref{cap:bool}), due volte annidata.}
Il \texttt{tt}/\texttt{ff} sui due argomenti \texttt{b1, b2} \`e zucchero per $\indB$ applicato due volte. Agda impone esaustivit\`a: i quattro casi $(\tto,\tto), (\tto,\ff), (\ff,\tto), (\ff,\ff)$ non sono una scelta, sono i quattro rami obbligatori della doppia eliminazione.

\paragraph{(b) Riduzione computazionale di \texttt{\&\&} (\S\ref{cap:and-comm}.2).}
Una volta che \texttt{b1} e \texttt{b2} sono costruttori concreti, le clausole di \texttt{\&\&} riducono i due lati del tipo identit\`a a forme normali. Questo accade \emph{durante il type checking}: il tipo che si deve abitare in ciascun caso \`e quello \emph{dopo} la riduzione.

\medskip

Vediamolo caso per caso. In ciascun caso il tipo da abitare diventa:

\begin{center}
\renewcommand{\arraystretch}{1.2}
\begin{tabular}{lll}
\toprule
\textit{Caso} & \textit{Tipo originale} & \textit{Dopo riduzione di \texttt{\&\&}} \\
\midrule
$(\tto,\tto)$ & $\tto \andd \tto \identita \tto \andd \tto$ & $\tto \identita \tto$ \\
$(\tto,\ff)$  & $\tto \andd \ff \identita \ff \andd \tto$   & $\ff \identita \ff$ \\
$(\ff,\tto)$  & $\ff \andd \tto \identita \tto \andd \ff$   & $\ff \identita \ff$ \\
$(\ff,\ff)$   & $\ff \andd \ff \identita \ff \andd \ff$     & $\ff \identita \ff$ \\
\bottomrule
\end{tabular}
\end{center}

\smallskip

($\tto \andd \ff$ riduce a $\ff$ per la prima clausola di \texttt{\&\&}; $\ff \andd \tto$ riduce a $\ff$ per la seconda; ecc.)

\paragraph{(c) Introduzione del tipo identit\`a (\S\ref{cap:identita}) via $\refl$.}
In ogni caso, dopo la riduzione, i due lati sono \textbf{definitoriamente uguali}. Quindi $\refl$ li abita (Regola~\ref{reg:ponte}, il ponte). Le quattro \texttt{refl} nel codice \emph{non} sono lo stesso \texttt{refl}: sono quattro istanze diverse della regola di introduzione di \S\ref{cap:identita}, applicate in contesti diversi.

\section{Perch\'e \emph{non} basta matchare solo \texttt{b1}}

Sembrerebbe pulito scrivere solo due clausole:

\begin{agda}
&&-comm tt b2 = ???
&&-comm ff b2 = ???
\end{agda}

Vediamo il caso \texttt{tt}. Il tipo da abitare \`e $\tto \andd b_2 \identita b_2 \andd \tto$.

\smallskip
\begin{itemize}
\item $\tto \andd b_2$ riduce a $b_2$ (clausola 1 di \texttt{\&\&}). \checkmark
\item $b_2 \andd \tto$ \textbf{non riduce}, perch\'e $b_2$ \`e una variabile. $\times$
\end{itemize}
\smallskip

Quindi il tipo, dopo tutte le riduzioni possibili, resta $b_2 \identita b_2 \andd \tto$. Per $\refl$ servirebbero due lati definitoriamente uguali. Non lo sono. La propriet\`a ``qualcosa \texttt{\&\&} \texttt{tt} $\identita$ qualcosa'' \`e \textbf{proposizionalmente} vera (esiste una prova) ma \textbf{non definitoriamente} vera (non riduce). $\refl$ non la abita.

\begin{oss}[Conclusione strutturale]
Per ``sbloccare'' la riduzione di $b_2 \andd \tto$, serve case-split anche su $b_2$. \textbf{Quattro casi sono strutturalmente obbligatori}, non un capriccio della sintassi.
\end{oss}

\section{Il quadro completo}

Cosa abbiamo fatto, in una riga: il tipo $(b_1\,b_2 : \mathbb{B}) \to b_1 \andd b_2 \identita b_2 \andd b_1$ \`e due $\Pi$ annidati (\S\ref{cap:pi}, \S\ref{cap:pi-azione}); per abitarlo servono $\lambda b_1 \to \lambda b_2 \to \text{prova}$; per costruire la prova in modo che $\refl$ (\S\ref{cap:identita}) basti, i due lati del $\identita$ devono essere definitoriamente uguali (\S\ref{cap:giudizi}); per farli ridurre, servono i costruttori concreti, ottenuti via eliminazione di $\mathbb{B}$ (\S\ref{cap:bool}) in entrambe le variabili; in ognuno dei quattro casi le regole di computazione di \texttt{\&\&} (\S\ref{cap:and-comm}.2) riducono i due lati a $\tto \identita \tto$ o $\ff \identita \ff$, e $\refl$ chiude.

\bigskip

\noindent\textbf{Tre F-I-E-C --- $\mathbb{B}$ (E), \texttt{\&\&} (C), $\_\identita\_$ (I) --- che si intrecciano in quattro righe.} Niente magia, niente trucchi: solo le quattro regole del libro.

\bigskip

\begin{oss}[Cosa cambier\`a su $\mathbb{N}$]
Le prove del cap.~1 di Stump cadono tutte in questa famiglia: case-split puro + riduzione + $\refl$. Su $\mathbb{N}$ (cap.~2 di Stump) entreranno in gioco \emph{due} cose nuove. Primo: il \emph{passo induttivo} vero --- il caso $\suc\,k$ richieder\`a di chiamare l'ipotesi induttiva esplicitamente, cio\`e la chiamata ricorsiva nella definizione della funzione che dimostra. L'eliminatore $\indN$ entrer\`a in scena con la sua ricorsione, non solo come case-split. Secondo: il \emph{matching su prove di uguaglianza} (\texttt{sym refl = refl} ecc., \S\ref{sec:notevoli-id}) e il cero~\ref{cero:pm} entreranno in scena, perch\'e su $\mathbb{N}$ molte propriet\`a richiedono $\mathsf{cong}$, $\mathsf{subst}$, riscritture. Tutto quello che abbiamo costruito qui rester\`a, e si aggiungeranno questi due livelli.
\end{oss}